\chapter{Perbandingan Senilai dan Data}\label{ch:g6:propdata}

Kalau tiga buku tulis berharga $6$ euro, maka enam buku tulis berharga
$12$: lipatduakan buku tulisnya, lipatduakan harganya. Besaran yang
berperilaku seperti ini disebut \emph{\omterm{def:g6:propdata:def}{sebanding}} --- salah satu
gagasan yang paling berguna dalam seluruh matematika, yang dikembangkan
lebih jauh pada \cref{ch:g7:prop}. Bab ini ditutup dengan membaca dan
menggambar diagram data sederhana.

\section{Besaran yang sebanding}

\begin{definition}[Perbandingan senilai]\label{def:g6:propdata:def}
Dua besaran disebut \emph{sebanding}\index{perbandingan senilai} jika
nilai besaran yang satu diperoleh dari nilai besaran yang lain dengan
mengalikannya selalu dengan bilangan yang \emph{sama}, yang disebut
\emph{koefisien perbandingan}.
\end{definition}

\begin{example}\label{ex:g6:propdata:table}
Buku tulis seharga $2$ euro masing-masing:
\begin{center}
\renewcommand{\arraystretch}{1.3}
\begin{tabular}{l|cccc}
buku tulis & $3$ & $5$ & $8$ & $12$ \\
\hline
harga (euro) & $6$ & $10$ & $16$ & $24$ \\
\end{tabular}
\end{center}
Setiap harganya adalah banyaknya buku tulis $\times 2$: koefisiennya
$2$ (harga satuannya). Cara cepat memeriksa bahwa sebuah tabel
\omterm{def:g6:propdata:def}{sebanding}: semua ``\omterm{def:g3:division:remainder}{hasil bagi} kolom'' $\frac{6}{3}, \frac{10}{5},
\frac{16}{8}, \frac{24}{12}$ harus sama --- di sini semuanya sama dengan $2$.
\end{example}

\begin{example}[Sebuah bukan-contoh]\label{ex:g6:propdata:nonexample}
Umur dan tinggi badan tidak \omterm{def:g6:propdata:def}{sebanding}: anak berumur $12$ tahun tidak
dua kali lebih tinggi daripada anak berumur $6$ tahun. Periksalah
dengan bilangan: $1.50$ m pada umur $12$ tahun dan $1.15$ m pada umur
$6$ tahun memberi \omterm{def:g3:division:remainder}{hasil bagi} $\frac{1.50}{12} = 0.125$ dan
$\frac{1.15}{6} \approx 0.19$ --- tidak sama.
\end{example}

\begin{method}[Melengkapi tabel perbandingan senilai]\label{met:g6:propdata:complete}
Tiga cara, yang bebas dipilih:
\begin{enumerate}
  \item \emph{koefisien}: carilah pengalinya dari satu kolom yang
        lengkap, lalu terapkan pada kolom yang lain;
  \item \emph{kolom}: sebuah kolom dapat dikalikan dengan bilangan,
        atau dua kolom dapat dijumlahkan, untuk menghasilkan kolom
        yang baru;
  \item \emph{kembali ke satuan}: carilah dulu nilai untuk
        \emph{satu} barang, lalu kalikan.
\end{enumerate}
\end{method}

\begin{example}\label{ex:g6:propdata:methods}
Lima cangkir yang sama berharga $15$ euro; berapa harga delapan cangkir?

\emph{Kembali ke satuan}: satu cangkir berharga $15 \div 5 = 3$ euro,
jadi delapan cangkir berharga $8 \times 3 = 24$ euro.

\emph{Kolom}: $8 = 5 + 3$; tiga cangkir berharga $\frac{3}{5}$ dari
$15$, yaitu $9$; jadi delapan cangkir berharga $15 + 9 = 24$ euro.
Jawaban yang sama, memang seharusnya begitu.
\end{example}

\section{Persen}

\begin{definition}[Persen]\label{def:g6:propdata:percent}
``$t\,\%$ dari suatu besaran'' berarti \omterm{def:g6:fractions:def}{pecahan} $\frac{t}{100}$ dari
besaran itu: $25\,\%$ dari $60$ adalah $\frac{25}{100} \times 60 = 15$.
Persen adalah \omterm{def:g6:propdata:def}{perbandingan senilai} dengan koefisien $\frac{t}{100}$.
\end{definition}

\begin{example}\label{ex:g6:propdata:percent}
Sebuah kelas berisi $30$ murid punya $40\,\%$ murid perempuan.
Banyaknya murid perempuan, selangkah demi selangkah: $40\,\%$ berarti
$\frac{40}{100} = 0.4$, dan $0.4 \times 30 = 12$ murid perempuan.
Patokan yang berguna: $50\,\%$ adalah setengah, $25\,\%$ \omterm{def:g3:division:half}{seperempat},
$10\,\%$ sepersepuluh --- jadi $10\,\%$ dari $30$ adalah $3$, dan
$40\,\% = 4 \times 10\,\%$ adalah $4 \times 3 = 12$: jawaban yang
sama, dihitung di kepala.
\end{example}

\section{Membaca dan menggambar diagram}

\begin{method}[Membaca tabel atau diagram]\label{met:g6:propdata:read}
Apa pun gambarnya --- tabel, diagram batang, atau grafik garis:
\begin{enumerate}
  \item bacalah judulnya dan satuan pada kedua sumbunya lebih dulu;
  \item untuk menjawab pertanyaan, carilah batang atau kolom yang
        tepat, lalu bacalah nilainya dengan cermat pada skalanya;
  \item untuk membandingkan, bandingkan tinggi batangnya --- tetapi
        periksalah bahwa sumbunya bermula dari $0$, kalau tidak
        gambarnya dapat menyesatkan.
\end{enumerate}
\end{method}

\begin{example}\label{ex:g6:propdata:chart}
Diagram batang di bawah ini menunjukkan banyaknya buku yang dipinjam
di perpustakaan sekolah selama satu minggu.

\begin{omfigure}
\begin{tikzpicture}
\begin{axis}[
  width=8.6cm, height=5.2cm,
  ybar, bar width=16pt,
  symbolic x coords={Sen,Sel,Rab,Kam,Jum},
  xtick=data,
  ymin=0, ymax=22,
  ytick={5,10,15,20},
  ylabel={buku yang dipinjam},
  tick label style={font=\small},
  label style={font=\small}]
  \addplot[fill=omDef!30, draw=omDef] coordinates
    {(Sen,12) (Sel,8) (Rab,17) (Kam,6) (Jum,15)};
\end{axis}
\end{tikzpicture}

{\small Diagram batang: satu batang untuk satu hari, tingginya sama dengan hitungannya.}
\end{omfigure}

Bacaan: hari yang paling sibuk adalah Rabu ($17$ buku). Selasa dan
Kamis bersama-sama mencakup $8 + 6 = 14$ buku --- lebih sedikit
daripada Jumat sendirian ($15$). Seluruhnya untuk satu minggu:
$12 + 8 + 17 + 6 + 15 = 58$ buku.
\end{example}

\begin{omfigure}
\begin{tikzpicture}
\begin{axis}[omaxis,
  width=8.2cm, height=5.4cm,
  xmin=0, xmax=9.4, ymin=0, ymax=28,
  xlabel={buku tulis}, ylabel={harga (euro)},
  xtick={1,...,9}, ytick={4,8,12,16,20,24}]
  \addplot[omDef, thick, domain=0:9] {3*x};
  \addplot[omDef, only marks, mark=*, mark size=1.5pt] coordinates
    {(1,3) (2,6) (3,9) (4,12) (5,15) (6,18) (7,21) (8,24)};
  \draw[dashed, omThm] (axis cs:6,0) -- (axis cs:6,18)
    -- (axis cs:0,18);
\end{axis}
\end{tikzpicture}

{\small Besaran yang \omterm{def:g6:propdata:def}{sebanding} punya grafik yang sangat mudah dikenali:
titiknya berbaris pada sebuah garis lurus yang melalui titik asal (di
sini buku tulis seharga $3$ euro masing-masing; garis putus-putusnya
membaca harga $6$ buku tulis: $18$ euro).}
\end{omfigure}

\section{Latihan}

\begin{exercise}[$\star$]\label{exo:g6:propdata:1}
Apakah tabelnya \omterm{def:g6:propdata:def}{sebanding}? Berilah alasan dengan menghitung \omterm{def:g3:division:remainder}{hasil baginya}.
\begin{center}
\renewcommand{\arraystretch}{1.2}
\begin{tabular}{l|ccc}
kg apel & $2$ & $3$ & $5$ \\
\hline
harga (euro) & $5$ & $7.5$ & $12.5$ \\
\end{tabular}
\qquad
\begin{tabular}{l|ccc}
umur (tahun) & $2$ & $4$ & $8$ \\
\hline
tinggi (cm) & $85$ & $103$ & $130$ \\
\end{tabular}
\end{center}
\end{exercise}

\begin{exercise}[$\star$]\label{exo:g6:propdata:2}
Empat croissant berharga $4.80$ euro. Carilah harga satu croissant,
lalu harga tujuh croissant.
\end{exercise}

\begin{exercise}[$\star$]\label{exo:g6:propdata:3}
Lengkapilah tabel \omterm{def:g6:propdata:def}{perbandingan senilai} ini (koefisiennya dulu!):
\begin{center}
\renewcommand{\arraystretch}{1.2}
\begin{tabular}{l|cccc}
liter bahan bakar & $10$ & $25$ & $?$ & $60$ \\
\hline
harga (euro) & $18$ & $?$ & $72$ & $?$ \\
\end{tabular}
\end{center}
\end{exercise}

\begin{exercise}[$\star$]\label{exo:g6:propdata:4}
Sebuah mobil memakai $6$ L bahan bakar setiap $100$ km. Berapa bahan
bakar untuk $250$ km? Untuk $350$ km? Sejauh apa ia dapat melaju dengan $27$ L?
\end{exercise}

\begin{exercise}[$\star$]\label{exo:g6:propdata:5}
Hitunglah di kepala, dengan memakai patokan pada \cref{ex:g6:propdata:percent}:
$50\,\%$ dari $84$; \quad $25\,\%$ dari $200$; \quad $10\,\%$ dari
$63$; \quad $30\,\%$ dari $70$.
\end{exercise}

\begin{exercise}[$\star$]\label{exo:g6:propdata:6}
Di sebuah sekolah berisi $450$ murid, $60\,\%$ makan di kantin. Berapa
murid itu? Berapa yang tidak?
\end{exercise}

\begin{exercise}[$\star$]\label{exo:g6:propdata:7}
Dengan diagram batang pada \cref{ex:g6:propdata:chart}: pada hari apa
saja buku yang dipinjam kurang dari $10$? Berapa buku lebih banyak
yang dipinjam pada hari Rabu daripada pada hari Kamis?
\end{exercise}

\begin{exercise}[$\star$]\label{exo:g6:propdata:8}
Suhu pada tengah hari dari Senin sampai Jumat adalah $14$, $16$, $13$,
$17$, $20$ derajat. Gambarlah diagram batang data ini (pilihlah skala
yang masuk akal), lalu bacalah hari yang paling panas.
\end{exercise}

\begin{exercise}[$\star\star$]\label{exo:g6:propdata:9}
Sebuah resep untuk $6$ orang memerlukan $450$ g tepung dan $3$ butir
telur. Sesuaikan untuk $10$ orang. (Untuk telurnya, berpikirlah dulu
sebelum menulis \omterm{def:g6:decimals:places}{bilangan desimal} butir telur!)
\end{exercise}

\begin{exercise}[$\star\star$]\label{exo:g6:propdata:10}
Dengan kelajuan tetap, seorang pengendara sepeda menempuh $24$ km dalam $1$ jam.
\begin{enumerate}
  \item Sejauh apa ia melaju dalam $2$ jam? Dalam setengah jam? Dalam
        $1$ jam $30$ menit?
  \item Berapa lama waktu yang ia perlukan untuk $60$ km?
\end{enumerate}
\end{exercise}

\begin{exercise}[$\star\star$]\label{exo:g6:propdata:11}
Sebuah toko menawarkan ``potongan $20\,\%$ untuk semuanya''. Hitunglah
potongannya dan harga barunya untuk: bola seharga $15$ euro; raket
seharga $40$ euro. Apakah harga barunya \omterm{def:g6:propdata:def}{sebanding} dengan harga
lamanya? Berapa koefisiennya?
\end{exercise}

\begin{exercise}[$\star\star\star$]\label{exo:g6:propdata:12}
Dua batang lilin yang sama tingginya dinyalakan pada saat yang sama.
Yang tebal habis terbakar dalam $6$ jam, yang tipis dalam $4$ jam,
masing-masing dengan lajunya sendiri yang tetap. Setelah berapa lama
lilin yang tipis tepat setengah tinggi lilin yang tebal? (Nyatakan
tinggi yang tersisa setelah $t$ jam sebagai \omterm{def:g6:fractions:def}{pecahan} dari tinggi awalnya.)
\end{exercise}

\section{Soal: Peta, denah, dan cara berbohong dengan diagram}

\begin{problem}[{Soal akhir pekan --- skala sebuah peta adalah
perbandingan senilai: panjang mengikuti koefisiennya, luas mengikuti
kuadratnya, dan diagram yang digambar buruk menipu mata}]
\label{pb:g6:propdata:1}
Setiap peta, denah lantai dan model menuruti satu aturan: panjang
yang sebenarnya dan panjang yang digambar itu \emph{\omterm{def:g6:propdata:def}{sebanding}}
(\cref{def:g6:propdata:def}). Koefisiennya punya nama yang terkenal
--- yaitu \emph{skala}\index{skala} --- dan ia menyembunyikan dua
perangkap yang sengaja dipasang soal ini: luas \emph{tidak} mengikuti
skalanya, dan persen, seperti diagram, sepenuhnya bergantung pada apa
yang menjadi pembandingnya.

\medskip
\noindent\textbf{Bagian I --- Membaca sebuah peta.}
Sebuah peta pendakian berskala $1:100\,000$: setiap sentimeter pada
peta mewakili $100\,000$ sentimeter di tanah.
\begin{enumerate}
  \item Ubahlah $100\,000$ cm menjadi kilometer. Berapa km yang
        diwakili $1$ cm peta ini?
  \item Dua desa terletak $7.5$ cm terpisah pada peta; tepi sebuah
        danau berupa kurva sepanjang $12$ cm. Sebutkan jarak yang
        sebenarnya.
  \item Sebuah jalan Romawi yang benar-benar lurus panjangnya $20$ km.
        Berapa panjangnya pada peta?
  \item Peta kedua mengumumkan ``$1$ cm untuk $5$ km''. Tulislah
        skalanya dalam bentuk $1: \,?$. Peta mana di antara keduanya
        yang menunjukkan lebih banyak kerincian untuk daerah yang sama?
  \item Buatlah tabel kecil (jarak pada peta $1$, $2$, $7.5$ cm
        melawan jarak yang sebenarnya) untuk peta pendakian itu, lalu
        jelaskan mengapa skala tepat merupakan \omterm{def:g6:propdata:def}{perbandingan senilai}
        dalam arti \cref{def:g6:propdata:def}. Berapa koefisien yang
        mengubah sentimeter-pada-peta menjadi kilometer?
\end{enumerate}

\medskip
\noindent\textbf{Bagian II --- Denah kamar tidur Zahra.}
Zahra menggambar denah kamar tidurnya --- persegi panjang $4$ m kali
$3$ m --- dengan skala $1:50$.
\begin{enumerate}[resume]
  \item Berapa ukuran kamarnya pada denah itu?
  \item Tempat tidurnya berukuran $2$ m kali $0.9$ m, dan pintunya
        selebar $80$ cm. Sebutkan ketiga ukuran itu pada denahnya.
  \item Sekarang perangkapnya. Hitunglah luas kamar yang sebenarnya
        dalam m$^2$, lalu luas denahnya dalam cm$^2$. Ubahlah luas
        yang sebenarnya menjadi cm$^2$
        (\cref{def:g6:measure:area}) lalu bagilah dengan luas
        denahnya: apakah \omterm{def:g3:division:remainder}{hasil baginya} $50$?
  \item Jelaskan bilangan yang kamu temukan: pada denah $1:50$, setiap
        sentimeter persegi kertas mewakili persegi sebenarnya
        berukuran $50$ cm kali $50$ cm. Berapa cm$^2$ sebenarnya itu?
        Nyatakan aturannya: ketika panjang dibagi $50$, luas
        dibagi\,\dots
  \item Sebuah permadani bundar menutupi $6$ cm$^2$ pada denahnya.
        Berapa luas yang sebenarnya ia tutupi, dalam cm$^2$ lalu
        dalam m$^2$?
\end{enumerate}

\medskip
\noindent\textbf{Bagian III --- Cara berbohong dengan diagram (dan
dengan persen).}
\begin{enumerate}[resume]
  \item Sebuah toko menjual senilai $105$ euro pada hari Sabtu dan
        $110$ pada hari Minggu. Manajernya menggambar diagram batang
        yang sumbu tegaknya bermula dari $100$: batang Sabtu naik $5$
        satuan kecil, batang Minggu $10$. Apa yang disarankan
        gambarnya tentang hari Minggu, dan apa yang sebenarnya
        terjadi? (Hitunglah kenaikan hari Minggu sebagai persen dari
        hari Sabtu, sampai persen terdekat.) Peringatan yang mana
        pada \cref{met:g6:propdata:read} yang diabaikan sang manajer?
  \item Gambarlah ulang (atau jelaskan) diagram yang jujur, dengan
        sumbunya bermula dari $0$: bagaimana perbandingan kedua
        batangnya sekarang?
  \item Sebuah koleksi bertambah dari $40$ menjadi $60$ perangko:
        hitunglah kenaikannya sebagai persen dari nilai awalnya. Lalu
        ia menyusut kembali dari $60$ menjadi $40$: hitunglah
        penurunannya sebagai persen dari nilai awal \emph{yang itu}.
        Mengapa kedua jawabannya berbeda, padahal $20$ perangkonya sama?
  \item Musim obral: sebuah mantel seharga $100$ euro mendapat
        ``potongan $20\,\%$'', dan di kasir ada tambahan ``potongan
        $10\,\%$ dari harga yang sudah dipotong''. Hitunglah harga
        akhirnya selangkah demi selangkah. Apakah potongan
        seluruhnya $30\,\%$? Jelaskan kepada pembelinya berapa
        sebenarnya.
  \item Penutupnya, kembali ke peta pada pertanyaan 4 ($1$ cm untuk
        $5$ km): sebuah hutan menempati $3$ cm$^2$ pada peta itu.
        Berapa luasnya yang sebenarnya, dalam km$^2$? Tutuplah dengan
        aturan seluruh soal ini: panjang mengikuti skalanya, luas
        mengikuti\,\dots

\end{enumerate}
\end{problem}
